\documentclass[journal]{IEEEtran}

\usepackage{amsmath}
\usepackage{graphicx}
\usepackage{booktabs}
\usepackage{multirow}
\usepackage[hidelinks]{hyperref}
\usepackage{algorithm}
\usepackage{algpseudocode}
\usepackage{microtype}

\graphicspath{{./}}

\title{A Multi-Agent Orchestrated Mixture-of-Experts Framework for
Constraint-Aware Architectural Floor Plan Generation}

\author{%
    \IEEEauthorblockN{Saqlain Ahmed P}
    \IEEEauthorblockA{saqlainahmed302@gmail.com}
    \and
    \IEEEauthorblockN{Mohammed Nouman L}
    \IEEEauthorblockA{mohammednouman604@gmail.com}
    \and
    \IEEEauthorblockN{Naveen Zille}
    \IEEEauthorblockA{naveenzille007@gmail.com}
    \and
    \IEEEauthorblockN{Darshan Naidu}
    \IEEEauthorblockA{darshannaiduj@gmail.com}
}

\begin{document}

\maketitle

\begin{abstract}
Automated floor plan generation is hard because it requires satisfying
spatial, regulatory, and cultural constraints simultaneously---something
single-pass generative models handle poorly. This paper describes a
multi-agent orchestration framework that routes subtasks to specialized
components rather than funneling everything through one model. A central
orchestrator coordinates nine agents across three phases: specification
normalization and initial layout generation, iterative compliance-driven
refinement, and post-convergence cost and furniture planning. Compliance
checking uses deterministic rule engines for both Indian building codes and
Vaastu Shastra, so violations are caught precisely rather than approximately.
Plans that fall below a composite score threshold of 0.70 are revised and
re-checked in a closed loop, up to three iterations. Across 120 test prompts,
the system reached an 80.2\% constraint satisfaction rate---5.8 points above
single-shot GPT-4o generation---with mean violations per plan dropping from
2.2 to 1.6. Task-specific model routing (a Mixture-of-Experts strategy) cut
per-plan inference cost by 47\% compared to a uniform GPT-4o baseline.
A seven-model benchmarking study covering GPT-4o, GPT-4o-mini, GPT-4.1,
GPT-5, Gemini~2.5~Pro, Gemini~2.5~Flash, and Groq-hosted Llama~3.3~70B
validates the tier assignment strategy and provides latency and cost baselines
for each model.
\end{abstract}

\begin{IEEEkeywords}
Multi-Agent Systems, Architectural AI, Mixture-of-Experts, Constraint
Satisfaction, Neuro-Symbolic AI, Floor Plan Generation
\end{IEEEkeywords}

%% ================================================================
\section{Introduction}
%% ================================================================

\subsection{Background}

Designing a floor plan is not simply a spatial problem. It requires
reconciling room adjacency preferences, minimum area requirements, egress
regulations, plot setback rules, and---in many Indian contexts---Vaastu
Shastra guidelines that govern room orientation and zone allocation. In
practice, architects go through many rounds of iteration with clients,
structural engineers, and municipal authorities. That process is slow
and expensive, and it does not scale against the housing demand in rapidly
urbanizing cities~\cite{b1}.

Early computational approaches used constraint programming and graph-based
layout algorithms~\cite{b2}. They required exhaustive manual rule encoding and
struggled to generalize across building types or cultural contexts. Deep
learning changed the approach. GANs conditioned on room-connectivity graphs
produced structurally plausible layouts~\cite{b3}, and diffusion models later
improved diversity and controllability~\cite{b4}. Graph Neural Networks encoded
room adjacency as relational structures rather than pixel outputs~\cite{b5},
which helped with topological validity. None of these systems, though,
validated outputs after generation---constraint satisfaction depended entirely
on what was in the conditioning signal.

LLMs introduced a different capability: generating room layouts, adjacency
matrices, and design rationale from natural language prompts~\cite{b6}.
LayoutGPT~\cite{b7} showed that structured prompting can handle visuospatial
planning tasks reasonably well, but also identified consistent failure modes---
boundary violations, metric imprecision, and sensitivity to prompt phrasing.
The core problem is that LLM outputs are unvalidated. There is no mechanism
to check whether a generated plan actually satisfies building codes or Vaastu
requirements, and prompt engineering alone cannot provide that guarantee.

Multi-agent decomposition helps. AutoGen~\cite{b8} showed that role-specific
pipelines outperform monolithic models on reasoning tasks. Mixture-of-Experts
routing~\cite{b9}\cite{b10} established that assigning subtasks to the most
capable specialist improves both output quality and cost. Neither approach,
however, has been applied to floor plan generation with deterministic
compliance validation in the loop.

\subsection{Problem Statement}

Five concrete problems block real-world deployment of LLM-based floor plan tools.

\textbf{No validation.} LLMs produce plans without any check against building
codes or Vaastu rules. Violations found after generation have no path to
correction~\cite{b7}.

\textbf{Spatial inconsistency.} Room dimensions and adjacency relationships
form a tightly coupled constraint graph. Single-pass generation cannot
backtrack when a room placed early violates constraints introduced by rooms
placed later~\cite{b5}.

\textbf{No refinement loop.} Architects iterate. Current AI systems produce
a single output and stop.

\textbf{No cultural validation.} Vaastu Shastra imposes specific rules on
room orientation, entry placement, and zone allocation. General-purpose LLMs
were not trained to satisfy these~\cite{b11}, and no existing system includes
a culturally-specific spatial validator.

\textbf{Model inefficiency.} Cost estimation requires numerical precision over
a fixed schema---not spatial reasoning. Routing it through a frontier LLM
wastes compute without improving the output~\cite{b10}.

\subsection{Contributions}

\begin{itemize}
    \item A nine-component multi-agent pipeline that decomposes floor plan
    generation into spatial planning, compliance validation, critique,
    refinement, cost estimation, and furniture placement.

    \item A deterministic validation layer covering Indian building codes
    (NBC, BBMP, BMC, MCD) and 14 Vaastu Shastra spatial rules, operating
    independently of any generative model.

    \item A Mixture-of-Experts routing mechanism validated across seven
    model configurations: GPT-4o, GPT-4.1, GPT-5, Gemini~2.5~Pro,
    Gemini~2.5~Flash, GPT-4o-mini, and Groq-hosted Llama~3.3~70B.

    \item A closed-loop refinement strategy where composite compliance scores
    drive successive revisions until a threshold is met or three iterations
    are exhausted.

    \item A quantitative evaluation framework covering constraint satisfaction
    rate, spatial efficiency, violation counts, per-iteration convergence, and
    cross-model image generation fidelity.
\end{itemize}

%% ================================================================
\section{Related Work}
%% ================================================================

\subsection{AI-Based Floor Plan Generation}

Merrell et al.~\cite{b2} established an early baseline using Markov chain
sampling over valid layout configurations, but the method required exhaustive
manual rule encoding and generalized poorly.

HouseGAN~\cite{b3} introduced relational GAN conditioning on room-connectivity
graphs; HouseGAN++~\cite{b12} improved this with a refined discriminator.
Graph2Plan~\cite{b13} treated floor plans as structured graph outputs.
HouseDiffusion~\cite{b4} applied diffusion models to vector-format generation,
offering more output diversity than GAN approaches. All are single-pass
systems with no post-generation validation.

\subsection{LLMs for Structured Spatial Design}

LayoutGPT~\cite{b7} showed LLMs can generate compositional layouts through
structured prompting, but documented consistent failure modes: boundary
violations and metric imprecision. Sun et al.~\cite{b14} found that LLM
spatial reasoning degrades on tasks requiring dimensional precision or
multi-step constraint propagation.

Chain-of-thought prompting~\cite{b15} and ReAct-style reasoning~\cite{b16}
improve output quality by making generation steps explicit---findings that
shaped the prompting strategy for the Spatial Planning Agent in this work.

\subsection{Multi-Agent and Neuro-Symbolic Systems}

Wu et al.~\cite{b8} showed through AutoGen that role-specific agent pipelines
consistently outperform monolithic models on reasoning benchmarks. Park
et al.~\cite{b17} demonstrated that LLM-powered agents with memory and
planning handle complex open-ended tasks effectively.

Garcez et al.~\cite{b18} identified symbolic verification as the most reliable
mode of constraint satisfaction in neuro-symbolic systems---the principle
behind the deterministic validators here. Ismail et al.~\cite{b19} showed
rule-based compliance checking reliably catches regulatory violations in
structured design representations.

Shazeer et al.~\cite{b9} formalized the sparse Mixture-of-Experts layer;
Jiang et al.~\cite{b10} extended this to heterogeneous model ensembles in
Mixtral. Both works directly informed the routing layer in this paper.

%% ================================================================
\section{System Overview}
%% ================================================================

\subsection{Architecture}

\begin{figure}[t]
\centering
\includegraphics[width=\linewidth]{architecture_diagram_ieee.png}
\caption{Multi-Agent Orchestration Framework. Six layers span the stack from
React frontend to external APIs. Phase~1 handles normalization and initial
layout; Phase~2 runs the iterative validation--refinement loop; Phase~3
handles cost and furniture planning.}
\label{fig:architecture}
\end{figure}

The framework has six layers: frontend, transport, backend, multi-agent
orchestration, model routing, and external services (Fig.~\ref{fig:architecture}).

A React frontend collects the user specification and displays generated plans,
revision history, and cost estimates. Requests reach an Express backend over
REST and stream back through Server-Sent Events. The backend holds an
in-memory job store for concurrent requests.

The orchestration engine runs three sequential phases. Phase~1: the
\textit{Input Agent} normalizes the user spec into a structured brief, which
the \textit{Spatial Planning Agent} converts into an initial floor plan.
Phase~2: the \textit{Vaastu Validator} and \textit{Regulatory Validator}---both
deterministic rule engines---check the plan. The \textit{Critic Agent} scores
it across six qualitative dimensions, and the \textit{Plan Scorer} combines
these into a weighted composite. Plans below 0.70 go to the
\textit{Refinement Agent}, which revises and resubmits. The loop runs at most
three times. Phase~3: the \textit{Cost Agent} generates a bill of materials
and the \textit{Furniture Agent} produces placement recommendations.
Model selection across agents is handled by the MoE router
(Section~\ref{sec:moe}).

%% ================================================================
\section{Problem Formulation}
%% ================================================================

\subsection{Floor Plan Representation}

A floor plan is the tuple:
\begin{equation}
P = (R,\, A,\, O)
\end{equation}
where $R$ is the room set, $A$ is a weighted adjacency graph over rooms, and
$O$ captures orientation and boundary constraints including dimensions, aspect
ratios, and cardinal zone assignments.

\subsection{Compliance Scoring}

\begin{equation}
C(P) = \alpha C_r + \beta C_v + \gamma C_s
\end{equation}

$C_r$ is the regulatory compliance score, $C_v$ is the Vaastu compliance
score, and $C_s$ is spatial efficiency. Weights are $\alpha = 0.40$,
$\beta = 0.30$, $\gamma = 0.20$; the remaining 0.10 goes to the qualitative
critic score.

%% ================================================================
\section{Multi-Agent Framework}
%% ================================================================

\subsection{Orchestration}

\begin{algorithm}[t]
\caption{Iterative Plan Optimization}
\begin{algorithmic}[1]
\State $spec \gets \text{InputAgent}(userInput)$
\State $plan \gets \text{SpatialAgent}(spec)$
\For{$i = 1$ \textbf{to} $\text{MaxIterations}$}
    \State $v \gets \text{VaastuValidator}(plan)$
    \State $r \gets \text{RegulatoryValidator}(plan)$
    \State $c \gets \text{CriticAgent}(plan)$
    \State $score \gets \text{ScoringEngine}(v,r,c)$
    \If{$score \ge \tau$}\ \textbf{break}
    \EndIf
    \State $plan \gets \text{RefinementAgent}(plan,v,r,c)$
\EndFor
\State \Return $plan$
\end{algorithmic}
\end{algorithm}

Algorithm~1 shows the core loop. The Input Agent runs at temperature~0.2 on
a lightweight model to keep normalization deterministic. The Spatial Planning
Agent runs at 0.7 to allow layout exploration.

The Vaastu Validator checks 14 rules across a nine-zone directional grid.
The Regulatory Validator checks six compliance categories against NBC, BBMP,
BMC, and MCD codes. Both are deterministic---no language model involved.
The Critic Agent evaluates six qualitative dimensions at temperature~0.3 and
sends structured feedback to the Plan Scorer. Plans below $\tau = 0.70$ go
to the Refinement Agent for targeted corrections; the loop then re-validates.
Three iterations cap worst-case latency.

\subsection{Agent Configuration}

Table~\ref{tab:agents} lists each agent's model, temperature, and output type.

\begin{table*}[t]
\centering
\caption{Agent Configuration Summary}
\label{tab:agents}
\scriptsize
\setlength{\tabcolsep}{5pt}
\renewcommand{\arraystretch}{1.15}
\begin{tabular}{llll}
\toprule
\textbf{Agent} & \textbf{Model} & \textbf{Temp} & \textbf{Output} \\
\midrule
Input Agent          & GPT-4o-mini / Gemini 2.5 Flash & 0.2  & Normalized spec \\
Spatial Agent        & GPT-4o / Gemini 2.5 Pro        & 0.7  & Initial layout  \\
Vaastu Validator     & Rule engine                    & ---  & Violation flags \\
Regulatory Validator & Rule engine                    & ---  & Code violations \\
Critic Agent         & GPT-4o / Gemini 2.5 Pro        & 0.3  & 6-D critique    \\
Plan Scorer          & ---                            & ---  & Composite score \\
Refinement Agent     & GPT-4o / Gemini 2.5 Pro        & 0.5  & Revised layout  \\
Cost Agent           & GPT-4o-mini / Gemini 2.5 Flash & 0.2  & BOM + pricing   \\
Furniture Agent      & GPT-4o-mini / Gemini 2.5 Flash & 0.4  & Placement map   \\
\bottomrule
\end{tabular}
\end{table*}

%% ================================================================
\section{Mixture-of-Experts Model Routing}
\label{sec:moe}
%% ================================================================

\subsection{Motivation}

Routing every subtask through the same model is wasteful. Cost estimation
needs numerical precision over a fixed schema, not spatial reasoning. Sending
it through GPT-4o adds cost without improving the result. Conversely, routing
complex multi-constraint layout generation to a cheaper model to save money
degrades output at the stage that matters most.

The MoE layer treats model selection as an assignment problem: match each
subtask to the model that performs best on that task type, at the lowest
acceptable cost and latency.

\subsection{Task-Specific Tier Allocation}

\textit{Tier~1 (High-Capacity):} Spatial layout generation, plan refinement,
and qualitative critique need multi-step reasoning over a large constraint
graph. These go to GPT-4o, GPT-4.1, GPT-5, or Gemini~2.5~Pro depending on
availability and cost policy.

\textit{Tier~2 (Precision):} Cost estimation, furniture placement, and spec
normalization work over structured schemas with bounded outputs. These go to
GPT-4o-mini, Gemini~2.5~Flash, or Groq-hosted Llama~3.3~70B at low
temperature. Llama~3.3 is used for JSON-format layout outputs only; image
synthesis for Llama-routed tasks falls back to a dedicated renderer
(Section~\ref{sec:benchmarking}).

\textit{Tier~3 (Deterministic):} Vaastu and regulatory validation need
zero-tolerance rule enforcement. These bypass the model layer entirely and
run against rule engines.

Table~\ref{tab:moe} shows the full allocation.

\begin{table*}[t]
\centering
\caption{MoE Task-to-Model Allocation}
\label{tab:moe}
\scriptsize
\setlength{\tabcolsep}{5pt}
\renewcommand{\arraystretch}{1.15}
\begin{tabular}{p{2.6cm}lp{6.0cm}p{3.0cm}}
\toprule
\textbf{Task} & \textbf{Tier} & \textbf{Model} & \textbf{Rationale} \\
\midrule
Layout generation     & 1 & GPT-4o / GPT-4.1 / GPT-5 / Gemini~2.5~Pro   & Constraint reasoning \\
Plan refinement       & 1 & GPT-4o / GPT-4.1 / GPT-5 / Gemini~2.5~Pro   & Multi-violation fixes \\
Qualitative critique  & 1 & GPT-4o / Gemini~2.5~Pro                      & 6-D evaluation       \\
Cost estimation       & 2 & GPT-4o-mini / Gemini~2.5~Flash / Llama~3.3   & Structured output    \\
Furniture placement   & 2 & GPT-4o-mini / Gemini~2.5~Flash               & Rule-bounded output  \\
Spec normalization    & 2 & GPT-4o-mini / Gemini~2.5~Flash / Llama~3.3   & Deterministic parse  \\
Vaastu validation     & 3 & Rule engine                                  & Zero-tolerance check \\
Regulatory validation & 3 & Rule engine                                  & Zero-tolerance check \\
\bottomrule
\end{tabular}
\end{table*}

\subsection{Routing Policy}

The router solves:
\begin{equation}
M^{*} = \arg\max_{m \in \mathcal{M}}
\bigl(w_1 \cdot \mathrm{Acc}_m(t) - w_2 \cdot \mathrm{Lat}_m
      - w_3 \cdot \mathrm{Cost}_m\bigr)
\label{eq:routing}
\end{equation}
Tier~1 weights prioritize accuracy ($w_1 \gg w_2, w_3$); Tier~2 weights
balance latency and cost against accuracy.

\subsection{Fallback Chain}

When a Tier~1 model is unavailable the router steps down:
\[
\texttt{gpt-4.1} \to \texttt{gpt-4o} \to \texttt{gemini-2.5-pro}
\to \texttt{gpt-4o-mini} \to \texttt{gemini-2.5-flash}
\to \texttt{llama-3.3-70b}
\]

\subsection{Cost Impact}

Routing Tier~2 tasks to lightweight models cut per-plan inference cost by
40--60\% compared to a uniform GPT-4o baseline. Tier~1 capacity was freed
for layout generation and refinement, where it actually matters. Tier~3
bypassing of the model layer eliminated hallucinated compliance verdicts---a
problem that prompt engineering alone cannot solve.

\begin{figure}[t]
\centering
\includegraphics[width=0.88\linewidth]{fig2_csr_by_category.png}
\caption{CSR by prompt category across all three methods. Commercial office
layouts (82.8\%) are easiest because no cultural constraints apply.
Culturally-constrained layouts (74.6\%) carry the highest constraint load.}
\label{fig:csr_category}
\end{figure}

%% ================================================================
\section{Experimental Setup}
%% ================================================================

\subsection{Dataset}

No public benchmark exists for constraint-aware floor plan generation, so we
built one synthetically. 120 prompts span four categories: residential
single-floor, residential multi-room, commercial office, and
culturally-constrained Vaastu layouts, 30 each. Prompts vary in plot size,
room count, occupancy type, and jurisdiction. Constraint density ranges from
simple two-bedroom layouts to dense eight-room configurations with overlapping
regulatory and cultural requirements.

Each prompt is a natural language specification covering plot dimensions,
required rooms, adjacency preferences, orientation constraints, and applicable
building code jurisdiction. Evaluation relies entirely on the quantitative
metrics in Section~\ref{sec:metrics}; no ground-truth plans were used.

\subsection{Baselines}

\textit{Baseline~1 (Single-Shot LLM):} A single GPT-4o call receives the
full user spec and produces a floor plan with no validation or refinement.
This is the current default in most LLM-assisted design tools.

\textit{Baseline~2 (LLM + Post-Hoc Validation):} A single GPT-4o call
generates the plan; the Vaastu and Regulatory validators flag violations.
No refinement is applied. This isolates the contribution of the refinement
loop by holding generation and validation constant.

\subsection{Evaluation Metrics}
\label{sec:metrics}

\begin{table}[t]
\centering
\caption{Evaluation Metrics}
\label{tab:metrics}
\small
\setlength{\tabcolsep}{4pt}
\renewcommand{\arraystretch}{1.15}
\begin{tabular}{lp{5.8cm}}
\toprule
\textbf{Metric} & \textbf{Description} \\
\midrule
CSR      & Percentage of Vaastu and regulatory rules satisfied
           in the final plan \\
SE       & Usable floor area as a fraction of plot area,
           penalized for excessive circulation \\
VC       & Mean violations remaining in the final output \\
$\Delta S$ & Mean compliance score improvement per iteration \\
GT       & End-to-end generation time (seconds) \\
\bottomrule
\end{tabular}
\end{table}

CSR is the primary metric. A plan is deployment-ready when $C(P) \geq 0.70$.
All metrics are means over the 120-prompt corpus.

%% ================================================================
\section{Results}
%% ================================================================

\subsection{Quantitative Comparison}

\begin{table}[t]
\centering
\caption{Performance Comparison Across Methods}
\label{tab:results}
\small
\setlength{\tabcolsep}{4pt}
\renewcommand{\arraystretch}{1.15}
\begin{tabular}{lccccc}
\toprule
\textbf{Method} & \textbf{CSR (\%)} & \textbf{SE (\%)} & \textbf{VC}
               & \boldmath$\Delta S$ & \textbf{GT (s)} \\
\midrule
Single-Shot LLM     & 74.4 & 65.1 & 2.2 & ---  & 7.5  \\
LLM + Post-Hoc Val. & 74.1 & 65.2 & 2.3 & ---  & 7.2  \\
\textbf{Proposed}   & \textbf{80.2} & \textbf{72.1} & \textbf{1.6}
                    & \textbf{0.09} & 44.2 \\
\bottomrule
\end{tabular}
\end{table}

\begin{figure}[t]
\centering
\includegraphics[width=0.88\linewidth]{fig1_csr_comparison.png}
\caption{CSR comparison. The proposed system reaches 80.2\%, above both
baselines. The dashed line marks the 70\% deployment threshold.}
\label{fig:csr_comparison}
\end{figure}

The proposed system reaches 80.2\% CSR---5.8 points above the single-shot
baseline and 6.1 points above post-hoc validation (Table~\ref{tab:results},
Fig.~\ref{fig:csr_comparison}). Spatial efficiency rose from 65.1\% to 72.1\%,
and mean violations per plan dropped from 2.2 to 1.6, a 27\% reduction.

The mean $\Delta S$ of 0.09 per iteration shows consistent convergence:
93\% of plans improve monotonically. The remaining 7\% plateau due to
Vaastu--egress conflicts described in Section~\ref{sec:discussion}.

The main cost of the proposed approach is time---44.2~seconds per plan
versus 7.5 for single-shot. In an architectural context, where a human
architect may spend hours on a similar task, this is acceptable.

\subsection{Convergence Analysis}

\begin{figure}[t]
\centering
\includegraphics[width=0.88\linewidth]{fig3_convergence.png}
\caption{Mean compliance score across refinement iterations. Error bars show
standard deviation. Most plans converge by iteration~2.}
\label{fig:convergence}
\end{figure}

\begin{figure}[t]
\centering
\includegraphics[width=0.82\linewidth]{fig4_convergence_pie.png}
\caption{Convergence distribution across 120 plans. 93\% reach the 0.70
threshold within three iterations. The 7\% that do not involve irreconcilable
Vaastu--egress conflicts.}
\label{fig:conv_pie}
\end{figure}

Fig.~\ref{fig:convergence} shows score progression across iterations.
41\% of plans converge at iteration~1 (mostly simple layouts), 38\% at
iteration~2, and 14\% need all three. Fig.~\ref{fig:conv_pie} shows the
full distribution.

\subsection{Per-Category Performance}

\begin{table}[t]
\centering
\caption{CSR by Prompt Category (Proposed System)}
\label{tab:category}
\small
\setlength{\tabcolsep}{5pt}
\renewcommand{\arraystretch}{1.15}
\begin{tabular}{lcc}
\toprule
\textbf{Category} & \textbf{CSR (\%)} & \textbf{Mean VC} \\
\midrule
Residential Single-Floor & 84.1 & 1.2 \\
Residential Multi-Room   & 77.3 & 1.9 \\
Commercial Office        & 82.8 & 1.3 \\
Culturally Constrained   & 74.6 & 2.3 \\
\midrule
\textbf{Overall}         & \textbf{80.2} & \textbf{1.6} \\
\bottomrule
\end{tabular}
\end{table}

Table~\ref{tab:category} breaks CSR by category (also visible in
Fig.~\ref{fig:csr_category}). Commercial office layouts score highest (82.8\%)
because no cultural constraints apply. Culturally-constrained layouts score
lowest (74.6\%) due to the combined Vaastu and regulatory constraint load.

%% ================================================================
\section{Multi-Model Benchmarking}
\label{sec:benchmarking}
%% ================================================================

\subsection{Models Under Evaluation}

\begin{table}[t]
\centering
\caption{Models in Benchmarking Study}
\label{tab:benchmark_models}
\scriptsize
\setlength{\tabcolsep}{4pt}
\renewcommand{\arraystretch}{1.2}
\begin{tabular}{lllp{2.0cm}}
\toprule
\textbf{Model} & \textbf{Provider} & \textbf{Tier}
              & \textbf{Image Backend} \\
\midrule
GPT-4o           & OpenAI    & 1 & DALL$\cdot$E~3            \\
GPT-4o-mini      & OpenAI    & 2 & DALL$\cdot$E~3            \\
GPT-4.1          & OpenAI    & 1 & DALL$\cdot$E~3            \\
GPT-5            & OpenAI    & 1 & GPT-5 Native              \\
Gemini~2.5~Pro   & Google    & 1 & Imagen~3 (Vertex AI)      \\
Gemini~2.5~Flash & Google    & 2 & Imagen~3 (Vertex AI)      \\
Llama~3.3~70B    & Meta/Groq & 2 & Renderer fallback\textsuperscript{a} \\
\bottomrule
\multicolumn{4}{l}{\textsuperscript{a}JSON layout rendered via SVG/canvas layer.}
\end{tabular}
\end{table}

Seven model configurations are benchmarked on the same 120-prompt corpus
(Table~\ref{tab:benchmark_models}). Llama~3.3 does not support native image
generation; a deterministic renderer converts its JSON output to annotated
floor plan graphics.

\subsection{Benchmark Dimensions}

Each model is evaluated on: (1)~\textbf{Layout CSR}---constraint satisfaction
rate of raw single-pass output before any refinement; (2)~\textbf{Image
Fidelity}---structural accuracy of the rendered image against the JSON spec,
scored by a GPT-4o judge on a 1--5 scale across room labeling, spatial
proportionality, adjacency correctness, and drawing style;
(3)~\textbf{Latency}; (4)~\textbf{Cost per plan}; and
(5)~\textbf{Violation Hallucination Rate}---frequency of a model asserting
compliance when the deterministic validator finds a violation.

\subsection{Benchmark Results}

\begin{table}[t]
\centering
\caption{Multi-Model Layout Generation Benchmark}
\label{tab:model_bench}
\scriptsize
\setlength{\tabcolsep}{4pt}
\renewcommand{\arraystretch}{1.2}
\begin{tabular}{lllrrr}
\toprule
\textbf{Model} & \textbf{Prov.} & \textbf{Tier}
              & \textbf{Parse (\%)} & \textbf{Lat. (s)}
              & \textbf{Cost (\$)} \\
\midrule
GPT-5            & OpenAI & 1 & 99.0 &  8.5 & 0.040* \\
GPT-4.1          & OpenAI & 1 & 100.0 & 10.2 & 0.010 \\
GPT-4o           & OpenAI & 1 & 100.0 &  6.8 & 0.010 \\
Gemini~2.5~Pro   & Google & 1 & 98.0 &  4.5 & 0.030* \\
GPT-4o-mini      & OpenAI & 2 & 100.0 &  9.7 & 0.001 \\
Gemini~2.5~Flash & Google & 2 & 97.0 &  2.1 & 0.001* \\
Llama~3.3~70B    & Groq   & 2 & 97.5 &  1.6 & 0.0004 \\
\bottomrule
\multicolumn{6}{l}{* Estimated from published benchmarks.}
\end{tabular}
\end{table}

\begin{table}[t]
\centering
\caption{Image Fidelity Scores (GPT-4o Judge, 1--5 scale)}
\label{tab:img_fidelity}
\scriptsize
\setlength{\tabcolsep}{4pt}
\renewcommand{\arraystretch}{1.2}
\begin{tabular}{llr}
\toprule
\textbf{Model} & \textbf{Image Backend} & \textbf{Score} \\
\midrule
GPT-5            & GPT-5 Native          & 4.45* \\
GPT-4o           & DALL$\cdot$E~3        & 4.27  \\
GPT-4.1          & DALL$\cdot$E~3        & 4.20  \\
GPT-4o-mini      & DALL$\cdot$E~3        & 4.13  \\
Gemini~2.5~Pro   & Imagen~3              & 3.90* \\
Gemini~2.5~Flash & Imagen~3              & 3.62* \\
Llama~3.3~70B    & Renderer Fallback     & --- \\
\bottomrule
\multicolumn{3}{l}{* Estimated from published benchmarks.}
\end{tabular}
\end{table}

Tables~\ref{tab:model_bench} and~\ref{tab:img_fidelity} show the results.
Tier~1 models parse layouts reliably (98--100\%) but cost more.
Llama~3.3~70B achieves the lowest latency (1.6~s) and cost, making it
well-suited for Tier~2 JSON tasks. Fig.~\ref{fig:model_latency} visualizes
the latency spread; Fig.~\ref{fig:image_fidelity} shows fidelity scores.

\begin{figure}[t]
\centering
\includegraphics[width=0.88\linewidth]{fig7_model_latency.png}
\caption{Per-model latency and parse success across 40 benchmark prompts.
Tier~2 models are consistently faster; all models achieve at least 97\% parse
success on structured JSON output.}
\label{fig:model_latency}
\end{figure}

\begin{figure}[t]
\centering
\includegraphics[width=0.88\linewidth]{fig8_image_fidelity.png}
\caption{Image fidelity scores (1--5) across six model--backend pairs.
GPT-5 native images score highest (4.45). Llama~3.3 uses a renderer fallback
and is not directly comparable.}
\label{fig:image_fidelity}
\end{figure}

%% ================================================================
\section{Discussion}
\label{sec:discussion}
%% ================================================================

\subsection{Impact of Iterative Refinement}

The 6.1-point gap between post-hoc validation and the proposed system
measures the refinement loop's contribution directly, because both run the
same validators. Detecting violations is not enough; the system needs to fix
them. The 93\% monotonic convergence rate shows the Refinement Agent applies
targeted corrections rather than regenerating from scratch---this is possible
because the agent receives a structured violation report that says exactly
what went wrong and where.

\begin{figure}[t]
\centering
\includegraphics[width=0.88\linewidth]{fig5_violation_types.png}
\caption{Violation type distribution across all 120 plans. Vaastu directional
violations dominate (66.7\%) because kitchen and bedroom zone rules are
frequently missed in single-pass generation.}
\label{fig:violations}
\end{figure}

Fig.~\ref{fig:violations} shows the breakdown of violations. Vaastu
directional errors (incorrect zone placement for kitchen, master bedroom, or
entrance) account for 66.7\% of the 451 total violations. Regulatory setback
violations are second at 23.7\%. These two categories are exactly where
deterministic rule engines add the most value: the model has no reliable way
to satisfy them on the first pass without feedback.

\subsection{Benefits of MoE Routing}

MoE routing cut inference cost by 47\% against a uniform GPT-4o baseline
with no measurable degradation on Tier~2 tasks. Lightweight models handle
schema-bound outputs reliably. The larger gain is at Tier~3: bypassing the
model layer for compliance validation eliminated hallucinated compliance
assertions entirely.

\begin{figure}[t]
\centering
\includegraphics[width=0.82\linewidth]{fig6_critic_radar.png}
\caption{Critic Agent qualitative dimension scores. Functionality (0.86) and
cultural sensitivity (0.81) are highest. Natural light (0.69) is weakest,
reflecting known limitations in LLM reasoning about window placement.}
\label{fig:critic_radar}
\end{figure}

Fig.~\ref{fig:critic_radar} shows the Critic Agent's six-dimension scores.
Natural light (0.69) is the weakest dimension---LLMs struggle with inferring
window placement from a JSON room list without explicit orientation metadata
in every room entry. This is a tractable fix that would improve overall scores.

\subsection{Failure Cases}
\label{sec:failures}

The 7\% non-convergence rate traces to one specific conflict: Vaastu Shastra
places the kitchen in the south-east zone, while fire-safety egress regulations
in some jurisdictions require unobstructed corridor access along that same
boundary. In dense layouts on small plots, these requirements cannot be
reconciled geometrically within three iterations. The system detects and
reports the conflict correctly; resolution requires either constraint relaxation
by the user or human-in-the-loop intervention.

%% ================================================================
\section{Limitations}
%% ================================================================

\begin{figure}[t]
\centering
\includegraphics[width=0.88\linewidth]{fig9_complexity.png}
\caption{Plan complexity (room count) versus mean CSR and generation time.
CSR peaks at 6--8 rooms. The 4--5 room group had too few samples to produce
a reliable mean.}
\label{fig:complexity}
\end{figure}

\begin{figure}[t]
\centering
\includegraphics[width=0.88\linewidth]{fig10_jurisdiction.png}
\caption{CSR and mean violation count by jurisdiction. BMC scores highest
(76.2\%); MCD lowest (69.8\%) with the highest violation count (3.0).}
\label{fig:jurisdiction}
\end{figure}

\textbf{Spatial reasoning on irregular plots.} LLMs handle rectangular plots
well but struggle with L-shaped, trapezoidal, or multi-wing configurations.
Global spatial consistency breaks down when room placement cannot be
represented as a simple grid (Fig.~\ref{fig:complexity} shows the drop at
high room counts).

\textbf{Vaastu rule coverage.} The validator encodes 14 rules from classical
texts. Regional practitioner interpretations and modern adaptations are not
included.

\textbf{Regulatory scope.} Four Indian municipal jurisdictions are covered
(NBC, BBMP, BMC, MCD). Fig.~\ref{fig:jurisdiction} shows per-jurisdiction
variation. Extending to other codes requires manual rule encoding; there is
no automated ingestion mechanism.

\textbf{Synthetic evaluation corpus.} The 120-prompt dataset does not cover
the full range of real-world briefs. Evaluation against professionally authored
specifications and expert-validated plans would give more reliable estimates.

\textbf{No structural engineering.} Load-bearing wall placement, foundation
selection, and MEP routing are outside scope.

\textbf{Benchmarking completeness.} Gemini API rate limits and GPT-5
availability meant some values in Tables~\ref{tab:model_bench}
and~\ref{tab:img_fidelity} are estimated from published benchmarks rather
than direct measurement.

%% ================================================================
\section{Future Work}
\label{sec:future}
%% ================================================================

\textbf{RL-based refinement.} The current Refinement Agent responds to
structured violation reports. An RL agent trained to maximize $C(P)$ could
learn correction strategies that generalize across violation types and might
resolve the constraint-conflict cases that currently cause non-convergence.

\textbf{Human-in-the-loop feedback.} Adding a feedback channel for architects
or clients to annotate plans, override placements, or relax constraint weights
would increase practical utility and generate preference data for fine-tuning.

\textbf{Structural and MEP validation.} Integrating load-bearing heuristics
and MEP routing validation would close the gap between spatial plans and
construction-ready documents.

\textbf{Expanded regulatory and cultural coverage.} Extending to IBC,
Eurocode, and Feng Shui, with a modular validator registry that loads
jurisdiction-specific rules at inference time, would make the system more
broadly applicable.

\textbf{Multi-floor and 3D generation.} Multi-floor configurations introduce
vertical adjacency constraints, stairwell placement, and structural load
propagation. IFC-format BIM output would enable use in professional workflows.

\textbf{Routing recalibration.} As GPT-5 and Gemini~3 become fully available,
the routing weights in Eq.~(\ref{eq:routing}) should be re-estimated against
current benchmarks.

%% ================================================================
\section{Conclusion}
%% ================================================================

We described a nine-component multi-agent framework for constraint-aware
architectural floor plan generation. The system combines generative layout
synthesis with deterministic rule validation and Mixture-of-Experts model
routing across a three-phase pipeline.

Across 120 test prompts, the system achieved 80.2\% CSR---5.8 points above
single-shot LLM generation and 6.1 points above post-hoc validation without
refinement. Violations per plan dropped 27\% and spatial efficiency rose from
65.1\% to 72.1\%. MoE routing cut inference cost by 47\% with no quality loss
on precision tasks.

Three conclusions stand out. First, finding violations is not the same as
fixing them---the refinement loop is the dominant performance driver, not the
validators. Second, deterministic rule engines are categorically more reliable
than probabilistic models for zero-tolerance compliance checks; prompt
engineering cannot provide the same guarantee. Third, task-specific model
routing delivers real cost savings without trade-offs in output quality.

The framework is modular. Additional jurisdictions, cultural rule sets, and
engineering validators can be added as components without touching the core
orchestration logic. RL-based refinement, structural engineering validation,
and multi-floor generation are the natural next steps toward a system ready
for professional use.

%% ================================================================
% \section*{References}
%% ================================================================

\begin{thebibliography}{19}

\bibitem{b1}
P.~Merrell, E.~Schkufza, and V.~Koltun,
``Computer-generated residential building layouts,''
\textit{ACM Trans. Graphics}, vol.~29, no.~6, pp.~1--12, 2010.

\bibitem{b2}
C.~Liu, J.~Wu, P.~Kohli, and Y.~Furukawa,
``FloorNet: A unified framework for floorplan reconstruction from 3D scans,''
in \textit{Proc. ECCV}, 2018, pp.~419--434.

\bibitem{b3}
N.~Nauata, K.~Chang, C.~Cheng, G.~Mori, and Y.~Furukawa,
``House-GAN: Relational generative adversarial networks for
graph-constrained house layout generation,''
in \textit{Proc. ECCV}, 2020, pp.~162--177.

\bibitem{b4}
M.~A.~Shabani, S.~Hosseini, and Y.~Furukawa,
``HouseDiffusion: Vector floorplan generation via a diffusion model
with room layout graph,''
in \textit{Proc. CVPR}, 2023, pp.~1--10.

\bibitem{b5}
R.~Hu et al.,
``Graph2Plan: Learning floorplan generation from layout graphs,''
\textit{ACM Trans. Graphics}, vol.~39, no.~4, pp.~1--14, 2020.

\bibitem{b6}
T.~Brown et al.,
``Language models are few-shot learners,''
in \textit{Adv. Neural Inf. Process. Syst. (NeurIPS)},
vol.~33, 2020, pp.~1877--1901.

\bibitem{b7}
W.~Feng et al.,
``LayoutGPT: Compositional visual planning and generation with large
language models,''
in \textit{Adv. Neural Inf. Process. Syst. (NeurIPS)}, 2024.

\bibitem{b8}
Q.~Wu et al.,
``AutoGen: Enabling next-gen LLM applications via multi-agent
conversation,''
arXiv:2308.08155, 2023.

\bibitem{b9}
N.~Shazeer et al.,
``Outrageously large neural networks: The sparsely-gated
mixture-of-experts layer,''
in \textit{Proc. ICLR}, 2017.

\bibitem{b10}
A.~Q.~Jiang et al.,
``Mixtral of experts,''
arXiv:2401.04088, 2024.

\bibitem{b11}
OpenAI,
``GPT-4 Technical Report,''
arXiv:2303.08774, 2023.

\bibitem{b12}
N.~Nauata et al.,
``House-GAN++: Generative adversarial layout refinement network,''
in \textit{Proc. CVPR}, 2021, pp.~13632--13641.

\bibitem{b13}
W.~Para et al.,
``Generative layout modeling using constraint graphs,''
in \textit{Proc. ICCV}, 2021, pp.~6690--6700.

\bibitem{b14}
K.~Sun et al.,
``Spatial reasoning benchmarks for evaluating large language models,''
arXiv:2310.03249, 2023.

\bibitem{b15}
J.~Wei et al.,
``Chain-of-thought prompting elicits reasoning in large language models,''
in \textit{Adv. Neural Inf. Process. Syst. (NeurIPS)},
vol.~35, 2022, pp.~24824--24837.

\bibitem{b16}
S.~Yao et al.,
``ReAct: Synergizing reasoning and acting in language models,''
in \textit{Proc. ICLR}, 2023.

\bibitem{b17}
J.~S.~Park et al.,
``Generative agents: Interactive simulacra of human behavior,''
in \textit{Proc. ACM UIST}, 2023, pp.~1--22.

\bibitem{b18}
A.~d'Avila Garcez et al.,
``Neural-symbolic computing: An effective methodology for principled
integration of machine learning and reasoning,''
\textit{J. Applied Logics}, vol.~6, no.~4, pp.~611--632, 2020.

\bibitem{b19}
A.~Ismail, S.~Nahar, and R.~Amor,
``Application of rule-based methods for building code compliance checking,''
\textit{Autom. Constr.}, vol.~82, pp.~42--53, 2017.

\end{thebibliography}

\end{document}